\section*{Technical Validation}

This section describes how we validate the extracted methodological fields released in \webisSRCorpus{}. The validation assesses whether values retained after the verify-then-repair pipeline are supported by the source full text. Since these fields are available only for the subset of reviews with accessible full texts, the validation reports extraction coverage separately from extraction correctness.

The section is organized in four steps. First, it summarizes descriptive statistics for the extracted subset, including field availability, abstention patterns, and major imbalances. Second, it defines a statistically justified extraction-validation design for the parsed full-text subset. Third, it evaluates extracted values using manual review as the primary source of evidence, with staged support from stronger language models and LLM-based judging where this helps scale the assessment without replacing human inspection. Fourth, it uses targeted sanity checks on notable outliers and failure cases to show where parsing, grounding, or verification remain unreliable.

The validation is designed to make three aspects of data quality transparent. First, non-null extracted values should be factually supported by the review text from which they were derived. Second, abstention rates should show how often the pipeline withholds a value when no verified extraction is retained. Third, the main sources of uncertainty should be visible, including parsing errors, weak evidence spans, and difficult field types. Together, these checks allow readers to assess which extracted values are reliable, which fields are sparse, and where additional caution is needed.
